% questions.tex —— 1985 年·全国卷·文科 · 题目内容片段（无答案）
% 由 tools/generate_projects.py 依来源总账自 1985/1985_数学真题_全国文_试题.md 生成，勿手改（改源 MD 后重跑）


\gkpart{文科数学}

\begin{gktip}
\textbf{来源说明}：本题面依据固定源 PDF《1985 全国文科数学》单页逐题转录。固定源为一页完整题面，原卷页码为 26；该 PDF 不含答案页，下列答案和解析由本库独立推导，不声称代表官方发布版本。
\end{gktip}

\begin{question}
\gkti{1}如果正方体 $ABCD-A'B'C'D'$ 的棱长为 $a$，那么四面体 $A'-ABD$ 的体积是\mbox{（　　）}

\nopagebreak[4]
\[
\text{（A）}\ \frac{a^3}{2}
\qquad
\text{（B）}\ \frac{a^3}{3}
\qquad
\text{（C）}\ \frac{a^3}{4}
\qquad
\text{（D）}\ \frac{a^3}{6}
\]

\end{question}
\begin{question}
\gkti{2}$\tan x=1$ 是 $x=\dfrac54\pi$ 的\mbox{（　　）}

（A）必要条件　　（B）充分条件　　（C）充分必要条件　　（D）既不充分又不必要的条件

\end{question}
\begin{question}
\gkti{3}设集合

\nopagebreak[4]
\[
X=\{0,1,2,4,5,7\},\qquad
Y=\{1,3,6,8,9\},\qquad
Z=\{3,7,8\},
\]

\Needspace{5\baselineskip}
那么集合 $(X\cap Y)\cup Z$ 是\mbox{（　　）}

\nopagebreak[4]
\[
\text{（A）}\ \{0,1,2,6,8\}
\qquad
\text{（B）}\ \{3,7,8\}
\]

\nopagebreak[4]
\[
\text{（C）}\ \{1,3,7,8\}
\qquad
\text{（D）}\ \{1,3,6,7,8\}.
\]

\end{question}
\begin{question}
\gkti{4}在下面给出的函数中，哪一个函数既是区间 $\left(0,\dfrac\pi2\right)$ 上的增函数又是以 $\pi$ 为周期的偶函数\mbox{（　　）}

（A）$y=x^2\ (x\in\mathbb R)$

（B）$y=|\sin x|\ (x\in\mathbb R)$

（C）$y=\cos2x\ (x\in\mathbb R)$

（D）$y=e^{\sin2x}\ (x\in\mathbb R)$

\end{question}
\begin{question}
\gkti{5}用 $1,2,3,4,5$ 这五个数字，可以组成比 $20000$ 大，并且百位数不是数字 $3$ 的没有重复数字的五位数，共有\mbox{（　　）}

（A）96 个　　（B）78 个　　（C）72 个　　（D）64 个

\end{question}
\begin{question}
\gkti{6}求函数

\nopagebreak[4]
\[
y=\frac{\sqrt{4-x^2}}{x-1}
\]

的定义域。

\end{question}
\begin{question}
\gkti{7}求圆锥曲线

\nopagebreak[4]
\[
3x^2-y^2+6x+2y-1=0
\]

的离心率。

\end{question}
\begin{question}
\gkti{8}求函数

\nopagebreak[4]
\[
y=-x^2+4x-2
\]

在区间 $[0,3]$ 上的最大值和最小值。

\end{question}
\begin{question}
\gkti{9}设

\nopagebreak[4]
\[
(3x-1)^6=a_6x^6+a_5x^5+a_4x^4+a_3x^3+a_2x^2+a_1x+a_0,
\]

求 $a_6+a_5+a_4+a_3+a_2+a_1+a_0$ 的值。

\end{question}
\begin{question}
\gkti{10}设 $i$ 是虚数单位，求 $(1+i)^6$ 的值。

\end{question}
\begin{question}
\gkti{11}设

\nopagebreak[4]
\[
S_1=1^2,\qquad
S_2=1^2+2^2+1^2,\qquad
S_3=1^2+2^2+3^2+2^2+1^2,\quad\ldots,
\]

\nopagebreak[4]
\[
S_n=1^2+2^2+3^2+\cdots+n^2+\cdots+3^2+2^2+1^2.
\]

\Needspace{5\baselineskip}
用数学归纳法证明：公式

\nopagebreak[4]
\[
S_n=\frac{n(2n^2+1)}{3}
\]

对所有的正整数 $n$ 都成立。

\end{question}
\begin{question}
\gkti{12}证明三角恒等式：

\nopagebreak[4]
\[
2\sin^4x+\frac34\sin^22x+5\cos^4x-\cos3x\cos x
=2(1+\cos^2x).
\]

\end{question}
\begin{question}
\gkti{13}解方程：

\nopagebreak[4]
\[
\lg(3-x)-\lg(3+x)=\lg(1-x)-\lg(2x+1).
\]

\end{question}
\begin{question}
\gkti{14}解不等式：

\nopagebreak[4]
\[
\sqrt{2x+5}>x+1.
\]

\end{question}
\begin{question}
\gkti{15}已知三棱锥 $V-ABC$ 的三个侧面与底面所成的二面角都是 $\beta$，它的高是 $h$。求这个棱锥底面的内切圆半径。

\end{question}
\begin{question}
\gkti{16}已知圆

\nopagebreak[4]
\[
C:x^2+y^2+4x-12y+39=0
\]

\Needspace{5\baselineskip}
和直线

\nopagebreak[4]
\[
L:3x-4y+5=0.
\]

求圆 $C$ 关于直线 $L$ 的对称的圆的方程。

\end{question}
\begin{question}
\gkti{17}设首项为 $1$、公比为 $q$（$q>0$）的等比数列的前 $n$ 项之和为 $S_n$，又设

\nopagebreak[4]
\[
T_n=\frac{S_n}{S_{n+1}},\qquad n=1,2,\ldots
\]

\Needspace{5\baselineskip}
求

\nopagebreak[4]
\[
\lim_{n\to\infty}T_n.
\]

\end{question}